\section*{Cluster A Step 39: Clean-Separation Derivation Test}

\paragraph{Grade.}
Finite-carrier diagnostic. The computation tests whether clean separation can be derived from a standard confining-sector requirement. It does not derive content or values.

\paragraph{Carrier.}
The Step-35 residual is rederived as \(12\): \(8\) in \(\texttt{2|3}\) and \(4\) in \(\texttt{4}\).

\paragraph{Requirement.}
The standardness requirement says that a surviving confining sector carries fermionic charged matter only; confining-charged massive vector content is non-standard.

\paragraph{Result.}
The standardness requirement leaves \(8\) survivors, all in \(\texttt{2|3}\). Relaxing it to allow confining-charged massive vectors leaves all \(12\), so the condition is load-bearing. However, standardness implies clean separation and clean separation implies standardness on this carrier.

\paragraph{Verdict.}
\(\texttt{INDEPENDENT\_INTRODUCED\_CIRCULAR}\). The standardness condition is extensionally equivalent to clean separation on the finite carrier, so it does not establish a more fundamental derivation.

\paragraph{Boundary.}
Clean separation remains a qualified introduced condition. The Step-38 structure-level result remains useful but caveated on this axis.
